% Part: set-theory
% Chapter: story
% Section: cumulative-approach

\documentclass[../../../include/open-logic-section]{subfiles}

\begin{document}

\olfileid{sth}{story}{approach}
\olsection{সঞ্চয়ী-পর্যায়ক্রমিক পদ্ধতি}

সেটগুলিকে কীভাবে বিভিন্ন পর্যায়ে বিন্যস্ত করব, তার একটু
বিস্তারিত বিবরণ এই:
\begin{quote}
  সেটগুলি বিভিন্ন \emph{পর্যায়ে} গঠিত হয়। প্রত্যেক পর্যায়~$S$-এর
  ক্ষেত্রে কিছু পর্যায় $S$-এর \emph{আগে} থাকে। পর্যায়~$S$-এ, $S$-এর আগের
  পর্যায়গুলিতে গঠিত সেটগুলি দিয়ে তৈরি প্রত্যেক সংগ্রহকে একটি
  সেট হিসেবে গঠন করা হয়। পর্যায়গুলিতে গঠিত সেট ছাড়া আর
  কোনো সেট নেই। \citep[p.~323]{Shoenfield:AST}
\end{quote}
এটি \emph{সেটের সঞ্চয়ী-পর্যায়ক্রমিক ধারণা}-র একটি রূপরেখা।
\olref[sth][][]{part}-এ যে আনুষ্ঠানিক সেটতত্ত্ব দেব, তার
ভিত্তি হবে এই ধারণা।

বিষয়টি আরেকটু খুঁটিয়ে দেখা যাক। শোয়েনফিল্ডের বর্ণনা
অনুযায়ী, প্রত্যেক পর্যায়ে আগের পর্যায়গুলি থেকে পাওয়া
সেট দিয়ে আমরা নতুন সেট গঠন করি। তাই প্রথম পর্যায়, অর্থাৎ
পর্যায়~$0$-এ, কোনো \emph{আগের} পর্যায় নেই; সুতরাং আগের
পর্যায় থেকে পাওয়া সেটও তো নেইই।\footnote{একটি প্রথম পর্যায়
\emph{আছে} বলে কেন ধরব? \olref[z][story]{sec}-এ
\stagesord{}-এর পাদটীকা দেখো।} ফলে আমরা কেবল একটি সেট
গঠন করি: !!{element}s-বিহীন সেট~$\emptyset$। পর্যায়~$1$-এ
আগের পর্যায় থেকে ঠিক একটি সেট পাওয়া যায়, তাই একমাত্র নতুন
সেটটি~$\{\emptyset\}$। পর্যায়~$2$-এ আগের পর্যায়গুলি থেকে
দুটি সেট পাওয়া যায়, আর আমরা দুটি নতুন সেট
$\{\{\emptyset\}\}$ এবং $\{\emptyset, \{\emptyset\}\}$
গঠন করি। পর্যায়~$3$-এ আগের পর্যায়গুলি থেকে চারটি সেট
পাওয়া যায়, তাই আমরা বারোটি নতুন সেট গঠন করি\ldots।
এভাবে সেটের সঞ্চয়ী-পর্যায়ক্রমিক ছবিটি কিছুটা নিচের মতো
দেখাবে (সংখ্যাগুলি পর্যায় বোঝায়):
\begin{center}
  \begin{tikzpicture}[scale=0.6]
  \tikzset{cut_here/.style={densely dotted}}
  \draw (-4,6) -- (-1, 0)-- (2,6);
  \draw[cut_here] (-5,8)--(-4,6);
  \draw[cut_here] (2,6)--(3,8);
  \draw(-1.5, 1)--(-0.5, 1);
  \draw(-2, 2)--(0, 2);
  \draw(-2.5, 3)--(0.5,3);
  \draw(-3, 4)--(1, 4);
  \draw(-3.5, 5)--(1.5, 5);
  \draw(-4, 6)--(2, 6);
  \node[label] at (0, 0) {\small 0};
  \node[label] at (0.5, 1) {\small 1};
  \node[label] at (1, 2) {\small 2};
  \node[label] at (1.5, 3) {\small 3};
  \node[label] at (2, 4) {\small 4};
  \node[label] at (2.5, 5) {\small 5};
  \node[label] at (3, 6) {\small 6};
  \end{tikzpicture}
\end{center}
তাহলে এই কাহিনি গ্রহণ করব কেন?

একটি কারণ, কাহিনিটি সহজবোধ্য এবং তা নিয়ে কাজ করা যায়।
সবচেয়ে আপাত ধারণা, অর্থাৎ অনানুষ্ঠানিক ধর্মনির্দেশে
সেট-গঠন ব্যর্থ হওয়ার পরে, এমন একটি ভালো বিকল্প আমাদের দরকার।

কিন্তু কাহিনিটি \emph{শুধু} সহজবোধ্য নয়। এই কাহিনির উপর
ভিত্তি করে তৈরি সেটতত্ত্ব \emph{সঙ্গত} হবে বলে বিশ্বাস করার
ভালো কারণও আছে। কারণটি এই।

সেটের সঞ্চয়ী-পর্যায়ক্রমিক ধারণা অনুযায়ী আমরা বিভিন্ন
পর্যায়ে সেট গঠন করি; আর তাদের !!{element}s হতে হবে
\emph{ইতিমধ্যেই} উপলব্ধ বস্তু। তাই যেকোনো পর্যায়~$S$-এর
জন্য সেটটি গঠন করতে পারি:
\[
  R_S = \Setabs{x}{x \notin x \text{ এবং $x$ $S$-এর আগে উপলব্ধ ছিল}}
\]
এখন রাসেলের কূটাভাসের প্রমাণের যুক্তি দেখায়, $R_S$ নিজে
পর্যায়~$S$-এর আগে উপলব্ধ ছিল না। এতে কোনো স্ববিরোধ নেই।
তার উপর, সেটের সঞ্চয়ী-পর্যায়ক্রমিক ধারণা গ্রহণ করলে
রাসেলের সেটটি গঠন করা যাবে—এমন \emph{প্রত্যাশাই} আমাদের
থাকা উচিত নয়। কারণ সেটি হবে এমন সব অ-স্বসদস্য সেটের
সেট, যেগুলি ``কখনো-না-কখনো উপলব্ধ হবে''। সংক্ষেপে,
এই ধারণা অনুযায়ী রাসেলের সেট গঠন করা যায় না—এটি
\emph{বিস্ময়কর} নয়; বরং আমরা \emph{এই ফলই আশা করব}।

%In one sense, then, the cumulative-iterative conception of set yields a response to Russell's Paradox which is rather like the \emph{predicativist}'s response. After all, the predicativist said that the collection of all non-self-membered sets$_n$ is a set$_{n+1}$, and this set$_n$ will not be self-membered. (Indeed, most predicativists will treat it as \emph{ungrammatical} to try to ask whether a set is self-membered.)
%
%But there is an important difference between the cumulative-iterative approach and the predicativist approach: the cumulative-iterative approach treats all of our entities as being \emph{of the same kind}. The predicativist will presumably have an empty set$_0$ (i.e., a set$_0$ with no !!{element}s), and an empty set$_1$ (i.e., a set$_1$ with no !!{element}s), and these will be \emph{different entities}. But on the cumulative-iterative approach, there is just one empty set, $\emptyset$, and it will be available at every stage of the hierarchy.

%Indeed, the Axiom of Extensionality, stated at the start of this chapter, will hold true of the elements in this hierarchy (so that there can be \emph{only one} ``empty'' set). But I do not intend to use the cumulative-iterative conception to justify Extensionality (see \cite{Boolos1971}). Again: I suggest that we should accept Extensionality, just because we are interested in extensional collections.

\end{document}
